\documentclass[11pt]{article}

% ----------------------------------------------------------------------
%  Research Paper Synthesis Graph -- shareable notes
% ----------------------------------------------------------------------
\usepackage[utf8]{inputenc}
\usepackage[T1]{fontenc}
\usepackage[margin=1in]{geometry}
\usepackage{lmodern}
\usepackage{microtype}
\usepackage{enumitem}
\usepackage{hyperref}
\usepackage{xcolor}

\definecolor{linkblue}{HTML}{1A5FB4}
\hypersetup{
  colorlinks=true,
  linkcolor=linkblue,
  urlcolor=linkblue,
  citecolor=linkblue,
  pdftitle={Research Paper Synthesis Graph},
  pdfauthor={},
}

\setlist{itemsep=2pt, topsep=4pt}

\title{\vspace{-2em}\textbf{Research Paper Synthesis Graph}}
\date{}
\author{}

\begin{document}
\maketitle
\vspace{-3em}

% ----------------------------------------------------------------------
\section{Why Graph RAG?}

Because the ``research problem'' is relational. If we ask a lexical /
vector RAG system \emph{``what is this paper about''}, it answers well.
But when we ask \emph{``what methods have been tried on problem $X$, and
which of them were limited by $Y$?''}~\cite{edge2024graphrag}, the answer
lives in the relationships \emph{between} papers, methods, and
limitations --- not in any single passage. A graph captures those edges
explicitly, so multi-hop and query-focused questions become traversals
rather than guesses.

% ----------------------------------------------------------------------
\section{Terms}

\begin{enumerate}[label=\arabic*.]
  \item \textbf{GROBID} --- \emph{GeneRation Of BIbliographic Data}. An
        open-source machine-learning tool that extracts and structures
        information from scholarly documents, primarily PDFs.
  \item \textbf{Cypher} --- a query language for graph databases (e.g.\
        Neo4j).
  \item \textbf{Kuzu} --- an embedded graph database. Runs in-process,
        stores the graph in a local directory, and speaks Cypher --- much
        like SQLite does for relational data.
  \item \textbf{KuzuGraphStore} --- the adapter that lets a RAG pipeline
        read from and write to a Kuzu graph.
\end{enumerate}

% ----------------------------------------------------------------------
\section{Pipeline (sketch)}

\begin{enumerate}[label=\arabic*.]
  \item \textbf{Ingest} --- parse PDFs with GROBID into structured
        metadata (title, authors, references, sections).
  \item \textbf{Build graph} --- create nodes (papers, authors, methods,
        problems) and edges (\texttt{CITES}, \texttt{USES\_METHOD},
        \texttt{LIMITED\_BY}) inside Kuzu.
  \item \textbf{Query} --- express research questions as Cypher
        traversals over the synthesis graph.
  \item \textbf{Synthesize} --- feed retrieved subgraphs to the LLM for
        query-focused summarization.
\end{enumerate}

% ----------------------------------------------------------------------
\section*{Citations}
\addcontentsline{toc}{section}{Citations}

\begin{thebibliography}{9}

\bibitem{edge2024graphrag}
Edge et al.,
\emph{``From Local to Global: A Graph RAG Approach to Query-Focused
Summarization''}.
Microsoft Research, 2024.
\href{https://arxiv.org/abs/2404.16130}{arXiv:2404.16130}.

\end{thebibliography}

\end{document}
